\addcontentsline{toc}{chapter}{Post Scriptum}
\chapter*{Post Scriptum}

\begin{flushright}
\small
\textit{``Vida é o que acontece com a gente, \\
enquanto estamos fazendo outros planos\dots''}
\\
-- Allen Saunders, citado por John Lennon\\        
na canção \textit{Beautiful Boy}.
\end{flushright}

Quando terminei de escrever esse pequeno punhado de lembranças da nossa família, acreditava que, em Cananéia, Paulo e eu encerraríamos nossa passagem por esse mundo mansa e dignamente como os velhos barcos de pesca que, vez ou outra, encontrávamos encalhados à beira dos canais.
Por já não terem condições de navegar, eram ali deixados por seus donos para se desfazerem nas mesmas águas que por longos anos venceram em duros embates e que agora, em terna homenagem, os embalavam docemente no vai-e-vem das marolas.
Ali ficavam, suas cores vivas lentamente desbotando sob a ação do sol e do sal, cumprindo seu destino na renovação do milagre da Criação: morrendo, serviam de abrigo e alimento para a vida fervilhante do mangue.

Cananéia me parecia o cenário perfeito para viver a velhice.
Pequenina, pobre e sossegada, distante das fúteis solicitações das cidades maiores, favorece o recolhimento, uma compreensão mais profunda da essência do viver e, principalmente, aninhada entre duas incomparáveis potestades da natureza, o oceano e a montanha, prontamente nos reduz às nossas reais proporções, experiência incontornável para quem começa a se dedicar ao inventário da própria existência.

Mas, não.
A Vida, mestra insuperada de todos os roteiristas, deve ter achado que esse era um fim muito óbvio.
Ou, quem sabe, talvez concordasse com os meus nunca esquecidos amigos jesuítas, que defendem só se tornar verdadeiramente humano aquele que vier a experimentar o avesso de tudo que escolheu para si.
O fato é que deixamos Cananéia, dezesseis anos depois, e viemos nos estabelecer em São Paulo.
Ainda que nos possibilite estar perto dos caçulas, os mais amorosos dos filhos, São Paulo parece ser a soma de tudo do que fugimos, do que tentamos evitar perambulando pelos confins desse país.
Essa cidade imensa nos priva da liberdade.
Ou será que é a velhice? Ou a soma dos dois? O fato é que aqui, mais do que em qualquer outro lugar por onde passamos, somos desafiados continuamente a lidar com os nossos limites.
Fuga possível, só para dentro.
O que, convenhamos, a esta altura, não deixa de ser conveniente.
Foi em razão desse recolhimento compulsório que decidi partilhar esse registro com vocês.
Também por sugestão da Paula, é verdade.
Concordamos em que ele pode ter alguma utilidade.


Sempre digo que nos momentos que mais exigiram de mim ao longo da vida, a consciência de que pulsa no meu sangue a vitalidade e a força dos nossos antepassados e a lembrança das suas palavras, pensamentos e atitudes é que me encorajaram a ir adiante.
Afinal, se eles conseguiram, por que não eu? Se esse punhado de relatos puder ter sobre vocês efeito semelhante, que bom seria!

Contudo, vou parar por aqui.
Antevejo ainda algumas batalhas pela frente, mas agora, não mais de natureza a serem compartilhadas, imagino.
Não porque envolvam segredos indignos, nada disso, mas porque a cada dia ficam mais pessoais e intransferíveis.
Como diz Madre Teresa em sua oração, é a etapa da existência que deixa de dizer respeito aos outros porque ``é entre você e Deus''.

A partir daqui, sigam vocês.
Recomendo.
É uma bela maneira de procurar um sentido para a vida.
Mas é, principalmente, um exercício de autoconhecimento capaz de revelar incríveis surpresas.
E de aprender aceitação.
Nunca resignação, jamais! ACEITAÇÃO, essa é a palavra.

\begin{flushright}
São Paulo, 28 de junho de 2025.
\end{flushright}